\documentclass[12pt]{article}
\usepackage[utf8]{inputenc}
\usepackage[T1]{fontenc}
\usepackage{hyperref}
\title{Nume Language Resources for K-12 Curriculum Development}
\author{Compiled by Claude Sonnet 5 (claude-sonnet-5) for Andr\'as Kornai}
\date{}
\begin{document}
\maketitle

\section{Introduction}
Nume (ISO 639-3 tgs; Glottocode nume1241) is an Oceanic language spoken by about 700 people on the northeast coast of Gaua island, Banks Islands, Vanuatu. It is also called Tarasag or Gog. It is distinct from its neighbours Mwerlap and Dorig.

\section{Primary Educational Resources}
No Nume primers, textbooks or school programmes were found.

\section{Linguistic and Reference Materials}
\begin{itemize}
\item Wikipedia on Nume, covering clusivity, number and directionals: \url{https://en.wikipedia.org/wiki/Nume_language}
\item Omniglot page on the Latin-script writing of Nume, with an alphabet chart and sample text: \url{https://www.omniglot.com/writing/nume.htm}
\item TLOPO database entry for Nume: \url{https://tlopo.clld.org/languages/nume}
\item Endangered Languages Project entry: \url{https://endangeredlanguages.com/lang/2427?hl=zh-cn}
\end{itemize}
No published Nume dictionary turned up in the searches. Alexandre Fran\c{c}ois's work on the languages of Gaua is cited in these sources, but I did not see a URL for it.

\section{Religious and Community Resources}
None were found.

\section{Digital and Audio Resources}
\begin{itemize}
\item Pangloss (CNRS) corpus of the languages of Gaua, with Nume recordings and traditional stories collected in 2003: \url{https://pangloss.cnrs.fr/corpus/Nume?lang=en}
\end{itemize}

\section{Curriculum Development Context}
The materials are thin: a few descriptive sources and an archived recorded corpus. A curriculum could use the Pangloss stories and the existing Latin orthography, but it would need new teaching materials made with the community.
\end{document}